\scanpage{38}
\mspage{37}
% Top: cancelled opening and tensor product.
\deleted{Soient $A$ et $B$ deux algèbres}

Le produit tensoriel gradué gauche $A\otimes B$ est un foncteur sur la catégorie $C$ \inserted{bivariante} des algèbres graduées. Ce produit tensoriel (malheureusement!) ne s'interprète pas comme un produit dans la catégorie bivariante: car on n'a pas d'hom de $C$ \unsure{canonique} $f:A\longrightarrow A\otimes A$, $f=(f^*,f_*)$, qui jouerait le rôle d'une application diagonale. On trouve néanmoins que c'est un produit \inserted{dans $C$} quand $A,B$ sont toutes deux des Alg. de Poincaré\ldots

% Middle: simplicial object and special condition.
Algèbre bivariante simpliciale: c'est un objet simplicial (pas \uline{$\frac12$} simplicial!) de la catégorie des Algèbres bivariantes, donc un foncteur contravariant de la catégor. des ens finis non vides dans $C$. Il est ``\uline{spécial}'' s'il transforme sommes disjointes en ``produits tensoriel'', i.e. si pour $I=I'\amalg I''$, les hom
\[
A(I')\xrightarrow{\alpha^*} A(I),\quad A(I'')\xrightarrow{\beta^*} A(I)
\]
[où $\alpha:I'\longrightarrow I$, $\beta:I''\longrightarrow I$) définissent un \uline{isom}
\[
A(I')\otimes A(I'')\xrightarrow{\sim}A(I'\amalg I'').
\]

% Bottom: Poincare restriction.
Si \deleted{$A(\Delta^0)$ est une Algèbre} on se borne à la sous-catégorie pleine des Alg. biv. simpliciales $A(\unsure{\bullet|\sigma})$ telles que $A(\unsure{\Delta^0})=A_0$ soit une Algèbre de Poincaré, on trouve une \uline{équivalence} de la catégorie des Alg. biv. simpliciales et des alg. de Poincaré.
